\documentclass[a4paper,12pt]{article}
\usepackage[utf8]{inputenc}
\usepackage[T1]{fontenc}
\usepackage{lmodern}
\usepackage{geometry}
\geometry{a4paper, margin=1in}
\usepackage{graphicx}
\usepackage{listings}
\usepackage{xcolor}
\usepackage{hyperref}
\usepackage{fancyhdr}
\usepackage{amsmath}

\pagestyle{fancy}
\fancyhf{}
\fancyhead[L]{Massimo Mantineo}
\fancyhead[C]{Hands-On 19}
\fancyhead[R]{A.A. 2024/2025}
\fancyfoot[C]{\thepage}

\title{Relazione Hands-On 19: Federated Learning con Ray.io e PyTorch}
\author{Massimo Mantineo (Matricola: 541924)}
\date{Aprile 2026}

\definecolor{codegreen}{rgb}{0,0.6,0}
\definecolor{codegray}{rgb}{0.5,0.5,0.5}
\definecolor{codepurple}{rgb}{0.58,0,0.82}
\definecolor{backcolour}{rgb}{0.95,0.95,0.92}

\lstdefinestyle{mystyle}{
    backgroundcolor=\color{backcolour},   
    commentstyle=\color{codegreen},
    keywordstyle=\color{magenta},
    numberstyle=\tiny\color{codegray},
    stringstyle=\color{codepurple},
    basicstyle=\ttfamily\footnotesize,
    breakatwhitespace=false,         
    breaklines=true,                 
    captionpos=b,                    
    keepspaces=true,                 
    numbers=left,                    
    numbersep=5pt,                  
    showspaces=false,                
    showstringspaces=false,
    showtabs=false,                  
    tabsize=2
}
\lstset{style=mystyle}

\begin{document}

\maketitle
\tableofcontents
\newpage

\section{Introduzione}
L'obiettivo di questo Hands-On è l'implementazione di un sistema di \textbf{Federated Learning (FL)} per la classificazione di immagini del dataset MNIST. A differenza del paradigma tradizionale centralizzato, il Federated Learning permette di addestrare modelli di Machine Learning su dati distribuiti, mantenendo la privacy dei dati locali. 

In questo esperimento, abbiamo utilizzato \textbf{Ray.io} come framework di orchestrazione distribuita e \textbf{PyTorch} per la definizione e l'addestramento del modello neurale.

\section{Stato dell'Arte e Teoria}
Il Federated Learning si basa sul concetto di portare il calcolo dai dati, invece dei dati al calcolo. Il processo tipico prevede:
\begin{enumerate}
    \item Un'entità centrale (\textbf{Aggregator}) che invia i pesi del modello globale ai partecipanti.
    \item I partecipanti (\textbf{Workers}) che addestrano localmente il modello sui propri dati.
    \item L'invio degli aggiornamenti (gradienti o pesi) all'Aggregator.
    \item L'aggregazione degli aggiornamenti tramite algoritmi come il \textbf{Federated Averaging (FedAvg)}.
\end{enumerate}

L'algoritmo FedAvg calcola la media pesata dei parametri del modello:
\begin{equation}
    w_{t+1} = \sum_{k=1}^{K} \frac{n_k}{n} w_{t+1}^k
\end{equation}
dove $n_k$ è il numero di esempi del worker $k$ e $n$ il totale degli esempi.

\section{Metodologia e Implementazione}
Il sistema è stato implementato con 1 Aggregator e 6 Worker distribuiti come \textbf{Ray Actors}.

\subsection{Partizionamento dei Dati}
Le 60.000 immagini del set di training di MNIST sono state suddivise in modo casuale (\textbf{IID split}) tra i 6 worker, assegnando 10.000 immagini a ciascuno. Ogni worker gestisce il proprio \texttt{DataLoader} locale.

\subsection{Architettura Distribuita}
\begin{itemize}
    \item \textbf{FederatedWorker}: Riceve i pesi globali, esegue 1 epoca di training locale e restituisce i pesi aggiornati.
    \item \textbf{FederatedAggregator}: Coordina i round di comunicazione, esegue l'aggregazione FedAvg e valuta il modello globale su un set di test centrale (10.000 immagini).
\end{itemize}

\section{Risultati Sperimentali}
Il sistema è stato testato per 3 round di comunicazione. I log mostrano un rapido incremento dell'accuratezza globale:

\begin{verbatim}
>>> Initializing Federated Learning with 6 workers...
>>> Round 1 starting...
Round 1 - Global Test Accuracy: 92.48%
>>> Round 2 starting...
Round 2 - Global Test Accuracy: 95.21%
>>> Round 3 starting...
Round 3 - Global Test Accuracy: 96.16%
Training finished.
\end{verbatim}

L'accuratezza finale del \textbf{96.16\%} ottenuta in soli 3 round dimostra l'efficacia dell'approccio federato nella convergenza del modello MLP.

\section{Conclusioni}
L'implementazione ha dimostrato come Ray.io semplifichi drasticamente la gestione di attori distribuiti per compiti di intelligenza artificiale. Il Federated Learning si è rivelato un metodo robusto per l'addestramento distribuito, riuscendo a mantenere prestazioni paragonabili alla versione centralizzata pur mantenendo i dati frammentati su worker indipendenti.

\end{document}
